%!TeX root=../thesis.tex

\chapter{Introduction}
\label{cap:introducao}

Software engineering organizations increasingly strive for fast value delivery through
autonomous, empowered teams. Concepts such as Team Topologies and DevOps promote
decentralized ownership and rapid feedback loops. Yet architecture decisions---those
that constrain the structure and quality attributes of a system---often remain
centralized in a handful of experienced practitioners. When requirements evolve
continuously, as they do in scaled products, teams face a tension between speed and
structural integrity: either they perform extensive upfront design that slows delivery
or they allow the system to drift toward what Foote and Yoder termed a ``Big Ball of
Mud,'' sacrificing long-term sustainability for short-term velocity~\cite{footeYoder1999}.

Understanding why this tension is so persistent requires acknowledging the organizational
pressures that teams routinely face. Foote and Yoder identify several forces that push
teams toward expedient, unstructured growth rather than deliberate design: tight delivery
schedules that preclude architectural reflection, heterogeneous experience levels within
teams, the compounding complexity introduced by rapid change, and cost constraints that
make dedicated architectural effort difficult to justify in the short term. These pressures
are not pathological---they are endemic to competitive software development---and any
practical approach to managing architectural uncertainty must operate within them.

Software architecture itself has been defined in multiple ways, reflecting how the
discipline has matured. Perry and Wolf's foundational 1992 account introduced a tripartite
model in which architecture is constituted by elements (the processing, data, and
connecting components), form (the weighted relationships among them), and rationale (the
constraints that justify the chosen structure)~\cite{perryWolf1992}. This model
established that architecture is not merely a set of structural choices but an expression
of reasoning under constraints. Decades later, Richards and Ford offered a more
practitioner-oriented definition: software architecture encompasses the structure of a
system together with the characteristics the system must support, the principles and
decisions that govern its design, and the underlying design decisions
themselves~\cite{richardsFord2020}. Both perspectives converge on the view that
architectural decisions are inherently bound to uncertainty---uncertainty about which
quality characteristics will matter most, which structural choices will best satisfy
them, and how the constraints will evolve.

Architectural uncertainty is a natural consequence of operating in environments where
business requirements, technology options, and quality-attribute trade-offs are only
partially known at any point in time. Despite its pervasiveness, this uncertainty
usually stays implicit. Existing methods for deriving and managing software architecture
depend heavily on the tacit knowledge of experienced architects, offering limited
guidance to less-experienced team members~\cite{souza2019}. The systematic mapping by
Souza et al.\ analyzed 39 primary studies on architectural derivation methods and
concluded that every method relied on architects' personal experience for its core
decisions, with tool support being either absent or only partially effective. This
reliance creates a bottleneck: when architectural expertise is scarce, teams either
make uninformed decisions or wait for overloaded experts, both of which degrade
delivery flow.

\section{Narrowing the Scope}
\label{sec:scope}

This thesis does not attempt to address software architecture management in general.
Instead, it focuses specifically on \textbf{software engineering teams working on a
scaled software product} within a single real-world organization. The target problem
is the lack of explicit, tool-supported management of architecture-related
uncertainties that could enable more shared and decentralized decision-making within
and across these teams. By restricting scope to concrete teams and a tangible tool,
the work responds directly to feedback received during the author's qualification
examination, which recommended narrowing the original research agenda from broad
surveys of centralization and decentralization in software architecture to a focused
empirical contribution centered on a specific artifact.

\section{ArchHypo and Arch-Stage}
\label{sec:archhypo-archstage}

The ArchHypo technique, proposed by Silva et al., reframes software architecture as a
collection of testable hypotheses rather than a fixed set of design
decisions~\cite{silva2024tse,silva2024ieeesw}. Each architectural hypothesis captures
an uncertainty---for instance, whether a chosen database technology will meet
scalability demands---along with an assessment of its uncertainty level and potential
impact. Teams then formulate technical plans, such as architectural spikes, tracer
bullets, or triggers, to reduce uncertainty or defer decisions to a ``responsible
moment'' when more evidence is available.

ArchHypo was developed using a Design Science Research (DSR) methodology and grew out
of patterns for agile architecture that were initially disseminated at software
engineering conferences and through corporate training programs~\cite{silva2024ieeesw}.
The technique thus carries both a theoretical grounding in DSR and a practical lineage
in agile community practice.

ArchHypo has been evaluated in two primary empirical settings. The first evaluation,
reported in an IEEE Software article~\cite{silva2024ieeesw}, examined the Catch Solve
startup---a company offering web testing and monitoring services---over a ten-month
observation period. Four hypotheses were managed, addressing availability, reusability
of test templates, usability, and scalability of virtual machines. This constituted the
first documented application of ArchHypo in a real operational environment. The study
demonstrated that making architectural uncertainty explicit guided decisions on when to
act versus when to defer, but relied entirely on manual tracking.

The second, more comprehensive evaluation, reported in the IEEE Transactions on Software
Engineering~\cite{silva2024tse}, took place within Jetsoft, a Brazilian software
development company, and focused on PropSEP, a mission-critical budget proposal system
for the Planning and Management Secretariat of the State of São Paulo. The system
processed approved budget figures on the order of 317~billion~BRL (approximately
60~billion~USD), underscoring the high stakes of its architectural decisions. The
evaluation ran from May to August 2022, spanning six sprints of seven business days
each, and involved eight team members (P1--P8). Ten architectural hypotheses were
identified, spanning six quality attributes: performance, security, reliability,
flexibility, usability, and team productivity. This evaluation produced nine findings
(detailed in Chapter~\ref{cap:fundamentacao}) that collectively demonstrate ArchHypo's
value while surfacing the adoption barriers that motivate the present thesis.

Of particular relevance to the present work are two findings from the TSE
evaluation~\cite{silva2024tse}. Finding~8 identified that changes in the team's way of
working constituted a significant obstacle: mixing architectural and functional tasks,
difficulty estimating hypothesis-related work, and the perception that the extra effort
was not directly tied to product delivery all created friction. Finding~9, which was
unanimous across all eight participants, was that the technique was hard to learn. Team
members struggled to map architectural concerns to the hypothesis format, to define
appropriate technical plan actions, and to articulate the assumptions underlying a
hypothesis. After the initial workshop, no participant felt ready to apply ArchHypo
independently without ongoing support.

Arch-Stage (originally named HypoStage) was conceived as a direct response to this
tooling gap. Developed as a Backstage Internal Developer Portal plugin in an
undergraduate capstone project supervised by the author of this thesis, Arch-Stage
provides structured interfaces for creating and editing hypotheses, assessing
uncertainty and impact on Likert scales, defining technical plans, and visualizing the
portfolio of architectural
hypotheses~\cite{correa2025}.\footnote{The initial design and implementation of
Arch-Stage, up to commit \texttt{ceee509}, is the contribution of Pedro Henrique
Mariano Corrêa's capstone project. Subsequent refinements reflect the author's
empirical use with real teams.} By embedding these capabilities within an IDP, the
tool integrates hypothesis management into the daily engineering workflow rather than
keeping it as a separate, easily forgotten activity.

Neither the Catch Solve nor the Jetsoft evaluation involved any tool-mediated support
for ArchHypo; both relied exclusively on manual means such as spreadsheets, meetings,
and informal documentation. This thesis therefore represents a distinct contribution:
an empirical evaluation of ArchHypo in its first tool-supported deployment with real
teams.

\section{Objective}
\label{sec:objetivo}

The objective of this work is to empirically evaluate and refine the Arch-Stage tool,
a support tool for the ArchHypo technique, by applying it with diverse software
engineering teams in a real-world organization with a scaled software product,
assessing how it supports managing architectural uncertainty and informing its
improvement.

\section{Research Questions}
\label{sec:research-questions}

% TODO: refine and finalize research questions after deciding whether the primary
% focus is (a) interaction/usability with Arch-Stage and ArchHypo or
% (b) patterns of architectural uncertainties managed using the tool.

The following research questions guide this investigation:

\begin{description}
  \item[RQ1] How do software engineering teams use Arch-Stage to identify and manage
  architectural hypotheses in practice?
  % TODO: consider splitting into sub-questions on identification vs. management.

  \item[RQ2] What refinements to Arch-Stage are needed to better support teams with
  varying levels of architectural experience?
  % TODO: decide whether to scope this to usability, workflow fit, or feature gaps.

  \item[RQ3] What types of architectural uncertainties do teams surface when using
  Arch-Stage, and how do these evolve over time?
  % TODO: this RQ may be secondary depending on the chosen evaluation focus.
\end{description}

\section{Thesis Structure}
\label{sec:estrutura}

The remainder of this thesis is organized as follows. Chapter~\ref{cap:fundamentacao}
presents the background on software architecture uncertainty, hypotheses engineering,
ArchHypo, the Arch-Stage tool, and the gap in tool support for architectural
derivation. Chapter~\ref{cap:trabalhos-relacionados} provides a contextual synthesis
of related work, drawing on the author's earlier mapping and multivocal review
efforts. Chapter~\ref{cap:design-pesquisa} describes the research design, including
the empirical context, data sources, and analysis approach.
% TODO: add references to results and conclusion chapters when they are written.
